\section{Research Design}
\label{sec:methodology_research_design}

This study used a quantitative experimental research design to evaluate intrinsic-reward methods for reinforcement learning.
The experiment compared a proposed intrinsic-reward family, Gated Learning-Progress Exploration (GLPE), against standard and intrinsic-motivation baselines under a fixed Proximal Policy Optimization (PPO) training backbone.
The independent variable was the exploration method used during training.
The dependent variables were extrinsic episodic return, sample efficiency, threshold-reaching reliability, and wall-clock efficiency.
All reported performance measurements were based on environment reward only; intrinsic reward was used only as a training signal.

The design consisted of three stages.
First, the intrinsic reward was specified as a transition-level signal that combines feature-space impact with region-local learning progress.
Second, all methods were trained on the same benchmark environments using matched PPO hyperparameters, environment-step budgets, and random seeds within each environment.
Third, saved checkpoints were evaluated offline with deterministic action selection to separate policy evaluation from training-time stochasticity.

The study followed a controlled comparison protocol.
For a fixed environment, the PPO architecture, value architecture, optimizer settings, discount factor, generalized advantage estimation parameter, and evaluation procedure were held constant across methods.
The intrinsic-reward coefficient and clipping range were also held constant across intrinsic methods within the same environment.
This control was necessary so that observed differences could be attributed primarily to the exploration signal rather than to changes in the underlying policy optimizer.

The proposed method was evaluated in two principal forms.
The gated form, GLPE, suppresses intrinsic rewards in latent regions that exhibit high prediction error but low learning progress.
The ungated form, GLPE (no gate), uses the same base intrinsic score without region suppression.
Additional ablations removed either the feature-space impact component or the learning-progress component, allowing the contribution of each term to be isolated.
A computational diagnostic variant was also used to measure the cost of repeatedly computing robust global medians for the gate.

The methodology emphasized reproducibility through fixed training budgets, explicit seed sets, deterministic execution where supported, offline checkpoint evaluation, and aggregation across independent runs.
Because reinforcement-learning outcomes can vary substantially across seeds, no single training trajectory was treated as sufficient evidence.
Instead, conclusions were based on curves and summary statistics aggregated over all completed seeds for each method--environment pair.